\documentclass[11pt,a4paper]{article}
\usepackage[utf8]{inputenc}
\usepackage[T1]{fontenc}
\usepackage[french]{babel}
\usepackage{geometry}
\geometry{margin=2.5cm}
\usepackage{hyperref}
\usepackage{listings}
\lstset{
  basicstyle=\ttfamily\small,
  breaklines=true,
  frame=single
}

\title{Partie REST --- Bibliothèque}
\date{}

\begin{document}
\maketitle

\section{Base URL et endpoints testés}
Le module REST expose ses routes sur :
\begin{lstlisting}
http://localhost:8081
\end{lstlisting}

Les appels à utiliser (Thunder Client / curl) sont listés dans le fichier :
\begin{lstlisting}
rest/request.json
\end{lstlisting}

\section{Comment les requêtes sont récupérées (pipeline)}
Pour chaque appel HTTP entrant, le chemin d'exécution est :
\begin{lstlisting}
HTTP -> Controller (@RestController) -> Service (@Service) -> Repository (Spring Data JPA) -> SQL (PostgreSQL)
\end{lstlisting}

\paragraph{Important.}
Dans ce document, \textbf{appel interne} = \textbf{appel Java} à un repository/service (méthode Spring Data),
car c'est ce qui structure concrètement le nombre d'opérations nécessaires pour fabriquer la réponse.

\section{Nombre exact d'appels internes par endpoint}
Les chiffres ci-dessous proviennent directement du code (appels repository effectués) :

\subsection{GET \texttt{/api/livres} (liste/recherche)}
\textbf{Objectif.} Récupérer des livres et éventuellement filtrer par \texttt{titre}, \texttt{genre}, \texttt{disponible}.

\textbf{Chemin code.} \texttt{LivreController.getLivres} $\rightarrow$ \texttt{LivreService.rechercher}.

\textbf{Appels internes exacts.}
\begin{itemize}
  \item \textbf{1 appel repository} : \texttt{livreRepository.findAll()}
  \item Les filtres \texttt{titre/genre/disponible} sont ensuite appliqués \textbf{en mémoire} (streams Java), donc \textbf{0 appel interne supplémentaire}.
\end{itemize}

\textbf{Conséquence.} Même avec des paramètres de filtre, le serveur charge d'abord \textbf{tous} les livres, puis filtre côté Java.

\subsection{GET \texttt{/api/livres/\{id\}} (détail)}
\textbf{Chemin code.} \texttt{LivreController.getLivre} $\rightarrow$ \texttt{LivreService.findById}.

\textbf{Appels internes exacts.}
\begin{itemize}
  \item \textbf{1 appel repository} : \texttt{livreRepository.findById(id)}
\end{itemize}

\subsection{POST \texttt{/api/livres/\{livreId\}/emprunter?adherentId=...}}
\textbf{Objectif.} Créer un emprunt, marquer le livre indisponible.

\textbf{Chemin code.} \texttt{LivreController.emprunterLivre} $\rightarrow$ \texttt{EmpruntService.emprunter}.

\textbf{Appels internes exacts.}
\begin{itemize}
  \item \textbf{1} : \texttt{livreRepository.findById(livreId)}
  \item \textbf{1} : \texttt{adherentRepository.findById(adherentId)}
  \item \textbf{(appel service)} : \texttt{livreService.markAsUnavailable(livreId)} qui fait :
    \begin{itemize}
      \item \textbf{1} : \texttt{livreRepository.findById(livreId)}
      \item \textbf{1} : \texttt{livreRepository.save(livre)}
    \end{itemize}
  \item \textbf{1} : \texttt{empruntRepository.save(emprunt)}
\end{itemize}

\textbf{Total exact} : \textbf{5 appels internes} (2 lectures + 1 lecture + 1 écriture + 1 écriture).

\subsection{PUT \texttt{/api/livres/emprunts/\{empruntId\}/retourner}}
\textbf{Objectif.} Clôturer un emprunt, marquer le livre disponible.

\textbf{Chemin code.} \texttt{LivreController.retournerLivre} $\rightarrow$ \texttt{EmpruntService.retourner}.

\textbf{Appels internes exacts.}
\begin{itemize}
  \item \textbf{1} : \texttt{empruntRepository.findById(empruntId)}
  \item \textbf{(appel service)} : \texttt{livreService.markAsAvailable(livreId)} qui fait :
    \begin{itemize}
      \item \textbf{1} : \texttt{livreRepository.findById(livreId)}
      \item \textbf{1} : \texttt{livreRepository.save(livre)}
    \end{itemize}
  \item \textbf{1} : \texttt{empruntRepository.save(emprunt)}
\end{itemize}

\textbf{Total exact} : \textbf{4 appels internes}.

\section{Nombre d'appels HTTP nécessaires pour ``composer'' une donnée (comparaison avec GraphQL)}
\subsection{Cas ``Query 1'' (liste filtrée)}
\textbf{REST} : \textbf{1 seul appel HTTP} (\texttt{GET /api/livres?genre=...&disponible=true}).

\subsection{Cas ``Query 2'' (livres + auteur + emprunts actifs)}
En GraphQL, tout est renvoyé en \textbf{1 requête}. En REST, avec les endpoints présents :
\begin{itemize}
  \item \textbf{1 appel HTTP} pour la liste : \texttt{GET /api/livres}
  \item \textbf{Les emprunts actifs} ne sont \textbf{pas exposés} par un endpoint REST de lecture dans ce module (il existe seulement \texttt{POST emprunter} et \texttt{PUT retourner}).
\end{itemize}
Donc, \textbf{la donnée complète de la Query 2 n'est pas reconstructible en REST} avec les routes actuelles.

\end{document}

